% ----------------------------------------------------------------
\subsection{Gaussian Model}
\label{sec:dyn_gaussian}
% ----------------------------------------------------------------

This section applies to the cells $(i,\ell)$ with \texttt{model} $=$
\texttt{gaussian}. Inside it, the subscript $i\ell$ is dropped from the per-cell
parameters and from the states; $k$ is the day.

\textbf{Model.} A mission increment is Gaussian,
$\mathcal{N}(\mu^{\inc}_{jk}, v^{\inc}_{jk})$ (inputs \texttt{mu} and \texttt{v});
the initial damage is $\mathcal{N}(\mu_0, v_0)$ and each replacement draw is
$\mathcal{N}(\mu^{\new}, v^{\new})$. A known damage has variance 0.
\textbf{Assumption N1 (independence).} The increments, the initial damage and the
replacement draws are mutually independent. A repair acts on the damage path by
the law of the cell (Section~\ref{sec:repair_laws}). The Gaussian law gives
negative damage a small probability; the input checks of
Section~\ref{sec:consistency_checks} ($\mu \geq 3\sqrt{v}$ for the initial
damage and for each increment) keep it at most $\Phi(-3) = 0.135\,\%$ for each of
these terms.

\textbf{Tracked states and repair laws.} The cell tracks the mean $\mu_k$ and the
variance $v_k$, and under $\ARDone$ their anchors $g^{\mu}_k$ and $g^{v}_k$. Both
laws are accepted. The exact day maps are
\begin{align}
  \mu_k &= \begin{cases}
    \mu_{k-1} + \delta\mu_k & \text{no intervention on day } k,\\
    (1-\rho)\,\mu_{k-1} & \text{repair, } \ARDinf,\\
    (1-\rho)\,\mu_{k-1} + \rho\, g^{\mu}_{k-1} & \text{repair, } \ARDone,\\
    \mu^{\new} & \text{replacement,}
  \end{cases}
  \label{eq:g_map_mu} \\
  v_k &= \begin{cases}
    v_{k-1} + \delta v_k & \text{no intervention on day } k,\\
    (1-\rho)^2\, v_{k-1} & \text{repair, } \ARDinf,\\
    (1-\rho)^2\, v_{k-1} + \bigl(1 - (1-\rho)^2\bigr)\, g^{v}_{k-1}
      & \text{repair, } \ARDone,\\
    v^{\new} & \text{replacement,}
  \end{cases}
  \label{eq:g_map_v}
\end{align}
with $\delta\mu_k = \sum_{j=1}^{M} x_{ijk}\,\mu^{\inc}_{jk}$ and
$\delta v_k = \sum_{j=1}^{M} x_{ijk}\,v^{\inc}_{jk}$.

\begin{proposition}[The state stays Gaussian]
\label{prop:g_exact}
Under N1, on every day of a schedule $D_k \sim \mathcal{N}(m_k, V_k)$, where
$(m_k, V_k)$ follow the maps~\eqref{eq:g_map_mu} and~\eqref{eq:g_map_v}, with
anchors equal to the mean and the variance just after the last intervention.
\end{proposition}

\begin{proof}
A mission adds an independent Gaussian term, and $\ARDinf$ scales the damage by
$1-\rho$; both keep the law Gaussian and give the first two cases. Under
$\ARDone$, $D^- = G + \Delta$, with $G$ the damage just after the last
intervention and $\Delta$ the increments added since; $G$ and $\Delta$ are
independent Gaussian variables by N1. Then $D^+ = G + (1-\rho)\Delta$ is Gaussian,
with mean $m_G + (1-\rho)(m^- - m_G)$ and variance
\[
  V_G + (1-\rho)^2\, V_\Delta = V_G + (1-\rho)^2\,(V^- - V_G)
  = (1-\rho)^2\, V^- + \bigl(1 - (1-\rho)^2\bigr)\, V_G .
\]
A replacement starts from the replacement draw. Therefore the law is Gaussian on
every day, with the moments of the maps.
\end{proof}

So $\ARDone$ is exact for this model. Its variance map uses the factor
$(1-\rho)^2$ and the variance anchor $g^{v}$. The rows are the generic
rows of Section~\ref{sec:encoding}, with the data below;
$V_{\max} = \tau^2/\kappa_G^2$, with $\kappa_G = \Phi^{-1}(1-\varepsilon)$.

\begingroup\small
\begin{longtable}{@{}llllll@{}}
\toprule
State & Day increment & Factor $a$ & Natural range & Initial value & Replacement value \\
\midrule
$\mu$; anchor $g^{\mu}$ & $\delta\mu_k$ & $1-\rho$ & $\tau$ & $\mu_0$ & $\mu^{\new}$ \\
$v$; anchor $g^{v}$ & $\delta v_k$ & $(1-\rho)^2$ & $V_{\max}$ & $v_0$ & $v^{\new}$ \\
\bottomrule
\end{longtable}
\endgroup

\subsubsection{Reliability row}

Gaussian cells need $\varepsilon < 0.5$, so that $\kappa_G = \Phi^{-1}(1-\varepsilon) > 0$.
For $D \sim \mathcal{N}(m, V)$,
$\Pr(D > \tau) = 1 - \Phi\bigl((\tau - m)/\sqrt{V}\bigr) \leq \varepsilon$ exactly
when $m + \kappa_G\sqrt{V} \leq \tau$. With $\kappa_G > 0$ and $\tau - \mu \geq 0$ this is
equivalent to the squared row, imposed on every day $k = 1, \ldots, T$:
\begin{equation}
  \kappa_G^2\, v_k \leq (\tau - \mu_k)^2, \qquad \mu_k \leq \tau.
  \label{eq:g_row}
\end{equation}

\emph{Validity.} Let $(m_k, V_k) \leq (\mu_k, v_k)$ be the true moments
(Proposition~\ref{prop:g_exact} and Section~\ref{sec:encoding}). If $V_k = 0$,
then $D_k = m_k \leq \tau$. Otherwise $\mu_k < \tau$, because $\mu_k = \tau$
forces $v_k = 0$, and
$\Pr(D_k > \tau) = 1 - \Phi\bigl((\tau - m_k)/\sqrt{V_k}\bigr)
\leq 1 - \Phi\bigl((\tau - \mu_k)/\sqrt{v_k}\bigr) \leq \varepsilon$, because
$(\tau - \mu)/\sqrt{v}$ decreases in $\mu$ and in $v$.

\textbf{Forms.} In the form \texttt{exact}, Eq.~\eqref{eq:g_row} is a
\emph{nonconvex} quadratic constraint: $(\mu, v) = (\tau, 0)$ and
$(0, V_{\max})$ satisfy it, and their midpoint $(\tau/2, V_{\max}/2)$ does not,
since $\tau^2/2 > \tau^2/4$. The solver handles it through Gurobi's nonconvex
quadratic support (parameter \texttt{NonConvex} $= 2$). The row has the form
$v_k \leq f(\mu_k)$ with the convex function
\[
  f(\mu) = \frac{(\tau - \mu)^2}{\kappa_G^2}, \qquad \text{side row } \mu \leq \tau,
\]
so the forms \texttt{tangent} and \texttt{pwl} of
Section~\ref{sec:linear_forms} apply, and both keep the side row.

\emph{Remark (the level $\varepsilon$ and the range $V_{\max}$).}
$V_{\max} = f(0) = \tau^2/\kappa_G^2$ is the largest variance that the row admits, at
$\mu = 0$, so it is the natural range of $v$ and the big-M constant of its rows;
the input checks of Section~\ref{sec:consistency_checks} keep $v_0$, its starting
anchor and $v^{\new}$ inside it.
$V_{\max}/\tau^2$ is $0.18$ at $\varepsilon = 0.01$, $0.61$ at $0.1$, $3.6$ at
$0.3$ and about $1600$ at $0.49$. It grows without bound as $\varepsilon \to 0.5$,
and at $\varepsilon \geq 0.5$ the squared row does not describe the reliability
requirement, because $\kappa_G \leq 0$. Hence the condition $\varepsilon < 0.5$
(Section~\ref{sec:consistency_checks}); the range is large only close to 0.5.

\textbf{Optional row.} If \texttt{v\_max\_user} is given for the cell, the row
$v_k \leq V^{\mathrm{user}}$ is added on every day.

\textbf{Facts used by the encoding and loop arguments.}
(M1) The day maps are nondecreasing in $(\mu, v, g^{\mu}, g^{v})$: their
coefficients lie in $[0, 1]$.
(M2) The true moments follow the maps exactly (Proposition~\ref{prop:g_exact}).
(M3) If the tracked states satisfy Eq.~\eqref{eq:g_row} and the true moments are
not larger, then $\Pr(D_k > \tau) \leq \varepsilon$ (the validity paragraph).
The loop row of Section~\ref{sec:loop} covers $\mu$, $v$ and, under $\ARDone$,
$g^{\mu}$ and $g^{v}$.

% ----------------------------------------------------------------
\subsection{Inverse Gaussian Model}
\label{sec:dyn_ig}
% ----------------------------------------------------------------

This section applies to the cells $(i,\ell)$ with \texttt{model} $=$
\texttt{inverse\_gaussian} (IG). Inside it, the subscript $i\ell$ is dropped from
the per-cell parameters and from the states; $k$ is the day.

\textbf{Model.} $\mathrm{IG}(\mu, \lambda)$ is the inverse Gaussian law with mean
$\mu > 0$ and shape $\lambda > 0$. Its density is
\[
  f(x; \mu, \lambda) = \sqrt{\frac{\lambda}{2\pi x^3}}\,
  \exp\!\Bigl(-\frac{\lambda\,(x - \mu)^2}{2\mu^2 x}\Bigr), \qquad x > 0,
\]
and its variance is $\mu^3/\lambda$. For $c > 0$,
$c\,\mathrm{IG}(\mu, \lambda) = \mathrm{IG}(c\mu, c\lambda)$. The cell has one
dispersion $\eta > 0$ (input \texttt{eta}), and a term of mean $a$ has the law
$\mathrm{IG}(a, a/\eta) = a\,Z$ with $Z \sim \mathrm{IG}(1, 1/\eta)$. Its density is
$\sqrt{a/(2\pi\eta x^3)}\,\exp\bigl(-(x - a)^2/(2\eta a x)\bigr)$, and $\eta$ is
its squared coefficient of variation. Every term has this law with its own mean:
a mission increment (mean $\mu^{\inc}_{jk}$), the initial damage ($\mu_0$) and
each replacement draw ($\mu^{\new}$). A known damage $a$ is the term $a \cdot 1$.
The terms may depend on each other in any way: this model needs no independence
assumption.

\textbf{Tracked state and repair laws.} The cell tracks the mean $\mu_k$, and
under $\ARDone$ its anchor $g^{\mu}_k$. Both laws are accepted. The exact day map
of the mean is
\begin{equation}
  \mu_k = \begin{cases}
    \mu_{k-1} + \delta\mu_k & \text{no intervention on day } k,\\
    (1-\rho)\,\mu_{k-1} & \text{repair, } \ARDinf,\\
    (1-\rho)\,\mu_{k-1} + \rho\, g^{\mu}_{k-1} & \text{repair, } \ARDone,\\
    \mu^{\new} & \text{replacement,}
  \end{cases}
  \label{eq:ig_map}
\end{equation}
with $\delta\mu_k = \sum_{j=1}^{M} x_{ijk}\,\mu^{\inc}_{jk}$. The rows are the
generic rows of Section~\ref{sec:encoding}, with the data below.

\begingroup\small
\begin{longtable}{@{}llllll@{}}
\toprule
State & Day increment & Factor $a$ & Natural range & Initial value & Replacement value \\
\midrule
$\mu$; anchor $g^{\mu}$ & $\delta\mu_k$ & $1-\rho$ & $\tau$ & $\mu_0$ & $\mu^{\new}$ \\
\bottomrule
\end{longtable}
\endgroup

\subsubsection{Reliability row}

Between replacements, the damage on day $k$ is
$D_k = \sum_t w_t X_t = \sum_t (w_t a_t)\, Z_t$, where $X_t = a_t Z_t$ are the
terms added since the last replacement (or day 0), and each $Z_t$ has the law of
$Z$ or equals 1. Each weight $w_t \in (0, 1]$ is a product of factors $1-\rho$:
$\ARDinf$ multiplies every weight by $1-\rho$, and $\ARDone$, with
$D^+ = G + (1-\rho)\,\Delta$, multiplies only the weights of the terms added since
the last intervention. When the starting anchor $g^{\mu}_0$ is below $\mu_0$, the
initial damage $\mu_0 Z_0$ is read as the two terms $g^{\mu}_0 Z_0$ and
$(\mu_0 - g^{\mu}_0) Z_0$. The mean $\sum_t w_t a_t$ is at most the tracked mean
$\mu_k$ (Section~\ref{sec:encoding}).

\begin{proposition}[Stop-loss order]
\label{prop:ig_order}
Let $D = \sum_t w_t a_t Z_t$ with repair weights $w_t \in (0, 1]$, means
$a_t \geq 0$, $\sum_t w_t a_t \leq \mu$, and each $Z_t$ either with the law of $Z$
or equal to 1, under any joint law. Then
$\mathbb{E}(D - d)_+ \leq \mathbb{E}(\mu Z - d)_+$ for every $d \geq 0$.
\end{proposition}

\begin{proof}
Let $c_t = w_t a_t$ and $A = \sum_t c_t$. If $A = 0$, then $D = 0$. Otherwise, $x \mapsto (x - d)_+$
is convex and nondecreasing, $Z \geq 0$, and
$(A - d)_+ \leq \mathbb{E}(AZ - d)_+$ by Jensen's inequality ($\mathbb{E}Z = 1$),
so
\[
  \mathbb{E}(D - d)_+
  = \mathbb{E}\Bigl(\sum_t \frac{c_t}{A}\,(A Z_t) - d\Bigr)_+
  \leq \sum_t \frac{c_t}{A}\,\mathbb{E}(A Z_t - d)_+
  \leq \mathbb{E}(A Z - d)_+
  \leq \mathbb{E}(\mu Z - d)_+ .
\]
\end{proof}

Let $q$ be the $(1-\varepsilon)$-quantile of $Z$, $\Pr(Z > q) = \varepsilon$,
and let $C = \mathbb{E}[Z \mid Z > q]$, the conditional value at risk (CVaR) of
$Z$ at the level $1 - \varepsilon$. The cap and the row, on every day
$k = 1, \ldots, T$, are
\begin{equation}
  \bar\mu = \frac{\tau}{C}, \qquad \mu_k \leq \bar\mu.
  \label{eq:ig_row}
\end{equation}

\begin{theorem}[CVaR cap]
\label{thm:ig_cap}
Let $D \geq 0$ satisfy $\mathbb{E}(D - d)_+ \leq \mathbb{E}(\mu Z - d)_+$ for every
$d \geq 0$, with $0 \leq \mu \leq \bar\mu$. Then $\Pr(D > \tau) \leq \varepsilon$.
\end{theorem}

\begin{proof}
If $\mu = 0$, then $\mathbb{E}D \leq 0$ and $D = 0$. Let $\mu > 0$ and $d = \mu q$.
Since $C > q$, $d < \mu C \leq \tau$. For every $x$,
$\mathbf{1}\{x > \tau\} \leq (x - d)_+/(\tau - d)$, and
$\mathbb{E}(Z - q)_+ = \varepsilon\,(C - q)$. Then
\[
  \Pr(D > \tau) \leq \frac{\mathbb{E}(D - d)_+}{\tau - d}
  \leq \frac{\mu\,\mathbb{E}(Z - q)_+}{\tau - \mu q}
  = \frac{\varepsilon\,\mu\,(C - q)}{\tau - \mu q} \leq \varepsilon,
\]
where the last step is $\mu C \leq \tau$.
\end{proof}

By Proposition~\ref{prop:ig_order} with $\mu = \mu_k$ and
Theorem~\ref{thm:ig_cap}, the row~\eqref{eq:ig_row} gives
$\Pr(D_k > \tau) \leq \varepsilon$ on every day, under both laws. The row is
linear: it is built in its exact form, and a value of \texttt{reliability\_impl}
at the cell is ignored (Section~\ref{sec:linear_forms}).

\textbf{Computation of $\bar\mu$} (offline, once per cell, before the model is
built). With $r(x) = 1/\sqrt{\eta x}$, the distribution function of $Z$ is
\[
  F(x) = \Phi\bigl(r(x)\,(x - 1)\bigr)
  + e^{2/\eta}\,\Phi\bigl(-r(x)\,(x + 1)\bigr).
\]
The second term overflows for small $\eta$; it is computed as
$\exp\bigl(2/\eta + \ln\Phi(-r(x)(x+1))\bigr)$, with $\ln\Phi$ from a log-space
routine (\texttt{scipy.special.log\_ndtr}).
\begin{enumerate}[leftmargin=1.5em, nosep]
  \item $q$ is the root of $1 - F(q) = \varepsilon$, found by bisection or by
        \texttt{brentq} on $\ln q$. The method \texttt{isf} of
        \texttt{scipy.stats.invgauss}, with \texttt{mu} $= \eta$ and
        \texttt{scale} $= 1/\eta$, gives the same value.
  \item The tail mean has the closed form
        $\mathbb{E}[Z;\, Z > q] = F(1/q)$, so
        $\bar\mu = \tau\,\varepsilon/F(1/q)$. Derivation: substituting into the
        density gives $z\, f_Z(z) = z^{-2} f_Z(1/z)$, and the change of variable
        $u = 1/z$ gives
        $\int_q^\infty z\, f_Z(z)\, dz = \int_0^{1/q} f_Z(u)\, du = F(1/q)$.
\end{enumerate}
Values of $\bar\mu/\tau$:
\begin{center}
\begin{tabular}{@{}lccccc@{}}
\toprule
 & $\eta = 0.05$ & $0.25$ & $1$ & $4$ & $10$ \\
\midrule
$\varepsilon = 0.001$ & 0.491 & 0.242 & 0.101 & 0.0370 & 0.0190 \\
$\varepsilon = 0.01$ & 0.570 & 0.322 & 0.156 & 0.0679 & 0.0401 \\
$\varepsilon = 0.1$ & 0.693 & 0.477 & 0.299 & 0.183 & 0.141 \\
\bottomrule
\end{tabular}
\end{center}

\emph{Remark (the largest cap for this order).} Let $\mu > \bar\mu$. For a tail
probability $\varepsilon' > \varepsilon$ close to $\varepsilon$, with
$(1-\varepsilon')$-quantile $q'$, still $\mu\,\mathbb{E}[Z \mid Z > q'] > \tau$.
The two-point law
$D = \mu\,\mathbb{E}\bigl[Z \mid \mathbf{1}\{Z > q'\}\bigr]$ is below $\mu Z$ in
the stop-loss order, by Jensen's inequality, and $\Pr(D > \tau) \geq \varepsilon'$.
So no larger cap is safe for every law that Proposition~\ref{prop:ig_order}
allows. That class also contains laws that no schedule produces, so the row can
reject some safe schedules.

\emph{Remark (the quantile cap is not the row).} The cap $\tau/q$ makes the model
law $\mu Z$ meet the target exactly: $\Pr(\mu Z > \tau) = \varepsilon$ at
$\mu = \tau/q$. A sum of terms below $\mu Z$ in the stop-loss order can have a
larger tail. Example: $\eta = 10$, $\varepsilon = 0.1$, so $\tau/q = 0.473\,\tau$;
three independent missions of equal mean, with total mean $\tau/q$, give
$\Pr(D > \tau) = 0.118$. The CVaR cap is $0.141\,\tau$ there. At
$\varepsilon = 0.01$ and $\eta$ from 0.05 to 10, $\bar\mu$ is 0.93 to 0.59 times
$\tau/q$.

\textbf{Facts used by the encoding and loop arguments.}
(M1) The day map~\eqref{eq:ig_map} is nondecreasing in $(\mu, g^{\mu})$.
(M2) The true mean follows it exactly.
(M3) If $\mu_k \leq \bar\mu$ and the true mean is not larger, then
$\Pr(D_k > \tau) \leq \varepsilon$ (Proposition~\ref{prop:ig_order},
Theorem~\ref{thm:ig_cap}). The loop row of Section~\ref{sec:loop} covers $\mu$
and, under $\ARDone$, $g^{\mu}$.
Each day of each repetition has its own set of terms; the argument above holds
for each of them, because it needs neither independence nor a fixed set of
terms.

% ----------------------------------------------------------------
\subsection{Wiener Model}
\label{sec:dyn_wiener}
% ----------------------------------------------------------------

This section applies to the cells $(i,\ell)$ with \texttt{model} $=$
\texttt{wiener}. Inside it, the subscript $i\ell$ is dropped from the per-cell
parameters and from the states; $k$ is the day.

\textbf{Model.} The damage is a Wiener process with drift, observed once a day. A
mission day adds a Gaussian increment $\mathcal{N}(\mu^{\inc}_{jk}, \sigma^2)$:
the mean comes from the input \texttt{mu}, and the diffusion constant
$\sigma > 0$ (input \texttt{sigma}) is the same for every mission, because it
describes the material of the component, not the route. The initial damage is
$\mathcal{N}(\mu_0, v_0)$ with $v_0 \geq 0$, and each replacement draw is
$\mathcal{N}(\mu^{\new}, v^{\new})$. \textbf{Assumption W1 (independence).} These
terms are mutually independent.

\textbf{Repairs act on the mean only} (a modelling choice). A repair shifts the
damage by a known amount: $D^+ = D^- - \rho\, m^-$ under $\ARDinf$ and
$D^+ = D^- - \rho\,(m^- - g^{\mu})$ under $\ARDone$, with $m^- = \mathbb{E}D^-$.
The law stays Gaussian, the variance is unchanged, and only a replacement lowers
it.

\textbf{Tracked states and repair laws.} The cell tracks $\mu_k$ and $v_k$, and
under $\ARDone$ the mean anchor $g^{\mu}_k$; the variance has no anchor. Both
laws are accepted. The mean follows the map~\eqref{eq:g_map_mu} of the Gaussian
model. The variance map is
\begin{equation}
  v_k = \begin{cases}
    v_{k-1} + \delta v_k & \text{no intervention on day } k,\\
    v_{k-1} & \text{repair, both laws},\\
    v^{\new} & \text{replacement,}
  \end{cases}
  \qquad \delta v_k = \sigma^2 \sum_{j=1}^{M} x_{ijk}.
  \label{eq:w_map_v}
\end{equation}
Under W1 the damage is $\mathcal{N}(m_k, V_k)$ on every day, with $(m_k, V_k)$
given by these maps: a mission adds an independent Gaussian term, and a repair or
a replacement acts as stated above.

The rows are the generic rows of Section~\ref{sec:encoding}, with the data below.
The variance takes the repair map $\psi = v_{k-1}$, that is the factor $a = 1$
under both laws, and no anchor. A depot day adds no increment, so for $v$ the
generic no-intervention and repair rows together equal the single row
\begin{equation}
  v_k \geq v_{k-1} + \delta v_k - V_{\max}\, x^r_k,
  \label{eq:w_vrow}
\end{equation}
which may replace them in the big-M form.

\begingroup\small
\begin{longtable}{@{}llllll@{}}
\toprule
State & Day increment & Factor $a$ & Natural range & Initial value & Replacement value \\
\midrule
$\mu$; anchor $g^{\mu}$ & $\delta\mu_k$ & $1-\rho$ & $\tau$ & $\mu_0$ & $\mu^{\new}$ \\
$v$; no anchor & $\delta v_k$ & $1$ & $V_{\max}$ & $v_0$ & $v^{\new}$ \\
\bottomrule
\end{longtable}
\endgroup

\emph{Remark (a repair of the path).} A repair of the damage path, as in the
Gaussian model, gives the same mean and the variance $(1-\rho)^2 V$ under
$\ARDinf$, or $(1-\rho)^2 V + (1 - (1-\rho)^2)\,V_G$ under $\ARDone$, with
$V_G \leq V$ the variance just after the last intervention. Both are at most $V$.
The reliability row increases in $v$, so the rows of this section are safe for
that reading as well.

\subsubsection{Reliability row}

The row is the Gaussian row~\eqref{eq:g_row}, with $\kappa_G = \Phi^{-1}(1-\varepsilon)$
and $\varepsilon < 0.5$, on every day $k = 1, \ldots, T$; its validity argument,
its forms (\texttt{exact}, \texttt{tangent}, \texttt{pwl}, with the same
function $f$ and side row) and the remark on $V_{\max} = \tau^2/\kappa_G^2$ are those of
Section~\ref{sec:dyn_gaussian}. If \texttt{v\_max\_user} is given for the cell,
the row $v_k \leq V^{\mathrm{user}}$ is added on every day.

\textbf{Facts used by the encoding and loop arguments.}
(M1) The day maps are nondecreasing in $(\mu, v, g^{\mu})$.
(M2) The true moments follow them exactly.
(M3) As for the Gaussian model. The loop row of Section~\ref{sec:loop} covers
$\mu$, $v$ and, under $\ARDone$, $g^{\mu}$.

\emph{Remark (the loop forces a replacement).} The variance grows by $\sigma^2$
on every mission day and falls only at a replacement. If the train flies $n$
missions in the operating phase and the component is not replaced there, then
$v_T \geq v_{H_1} + n\,\sigma^2 > v_{H_1}$, and the loop row $v_T \leq v_{H_1}$
fails. So every feasible schedule includes a replacement of the component in the
operating phase whenever its train flies a mission there. No input check enforces this. With
\texttt{allow\_replacement} false, the train of a Wiener cell cannot fly in the
operating phase.

% ----------------------------------------------------------------
\subsection{Gamma Model}
\label{sec:dyn_gamma}
% ----------------------------------------------------------------

This section applies to the cells $(i,\ell)$ with \texttt{model} $=$
\texttt{gamma}. Inside it, the subscript $i\ell$ is dropped from the per-cell
parameters and from the states; $k$ is the day.

\textbf{Model.} $\mathrm{Gamma}(a, \beta)$ is the Gamma law with shape $a > 0$
and rate $\beta > 0$; its mean is $a/\beta$. Every term that adds damage is a
Gamma variable with its own rate:
\begin{itemize}[leftmargin=1.5em, nosep]
  \item a mission increment with mean $\mu^{\inc}_{jk} > 0$ (input \texttt{mu})
        and rate $\beta_{jk}$ (inputs \texttt{gamma\_beta} and
        \texttt{gamma\_beta\_trans}) is
        $\mathrm{Gamma}(\beta_{jk}\,\mu^{\inc}_{jk},\, \beta_{jk})$;
  \item the initial damage with mean $\mu_0 > 0$ is
        $\mathrm{Gamma}(\beta_0\,\mu_0,\, \beta_0)$ when its rate $\beta_0$ is
        given (input \texttt{gamma\_beta\_0}); when \texttt{gamma\_beta\_0} is
        omitted, the initial damage is known: the constant $\mu_0$;
  \item each replacement draw with mean $\mu^{\new} > 0$ is
        $\mathrm{Gamma}(\beta^{\new}\mu^{\new},\, \beta^{\new})$ when its rate
        $\beta^{\new}$ is given (input \texttt{gamma\_beta\_new}); when
        \texttt{gamma\_beta\_new} is omitted, it is the known constant
        $\mu^{\new}$.
\end{itemize}
\textbf{Assumption G1 (independence).} All these terms are mutually independent.

A repair acts on the damage path by the law of the cell
(Section~\ref{sec:repair_laws}): $\ARDinf$ gives $D^+ = (1-\rho)\,D^-$, and
$\ARDone$ gives $D^+ = D^{\mathrm{last}} + (1-\rho)\,(D^- - D^{\mathrm{last}})$.
After a repair the damage is, in general, not a Gamma variable: it is a weighted
sum of independent Gamma terms and known constants. The reliability row below
holds for every such sum.

\textbf{Tracked state and repair laws.} The cell tracks the mean $\mu_k$ only,
and under $\ARDone$ its anchor $g^{\mu}_k$. Both laws are accepted. The shape of
the damage is not tracked. The exact day map of the mean is
\begin{equation}
  \mu_k = \begin{cases}
    \mu_{k-1} + \delta\mu_k & \text{no intervention on day } k,\\
    (1-\rho)\,\mu_{k-1} & \text{repair, } \ARDinf,\\
    (1-\rho)\,\mu_{k-1} + \rho\, g^{\mu}_{k-1} & \text{repair, } \ARDone,\\
    \mu^{\new} & \text{replacement,}
  \end{cases}
  \label{eq:gm_map}
\end{equation}
with $\delta\mu_k = \sum_{j=1}^{M} x_{ijk}\,\mu^{\inc}_{jk}$. The map is exact
for the true mean $\mathbb{E}D_k$, because expectation is linear. The rows are
the generic rows of Section~\ref{sec:encoding}, with the data below.

\begingroup\small
\begin{longtable}{@{}llllll@{}}
\toprule
State & Day increment & Factor $a$ & Natural range & Initial value & Replacement value \\
\midrule
$\mu$; anchor $g^{\mu}$ & $\delta\mu_k$ & $1-\rho$ & $\tau$ & $\mu_0$ & $\mu^{\new}$ \\
\bottomrule
\end{longtable}
\endgroup

\textbf{Smallest rate.} The cell's smallest rate is
\begin{equation}
  \beta^* = \min\Bigl(\{\beta_{jk} : \mu^{\inc}_{jk} > 0\}
    \cup \{\beta_0 : \mu_0 > 0\} \cup \{\beta^{\new} : \mu^{\new} > 0\}\Bigr),
  \label{eq:gm_betastar}
\end{equation}
where $j$ runs over the missions and $k$ over the days of both phases, and
$\beta_0$, $\beta^{\new}$ enter only when they are given. So $\beta^*$ is the
smallest rate among the terms with a positive mean: a mission with
$\mu^{\inc}_{jk} = 0$ adds no term, and an omitted rate declares a known damage,
which is a constant and does not enter $\beta^*$. If no term has a rate, the
damage is a known constant, and the row below is $\mu_k \leq \tau$. Write
$t = \beta^*\tau$. Units of $1/\beta^*$ turn every law below into a law of rate
at least 1.

\subsubsection{Class of true states}

\begin{definition}[Class of weighted Gamma sums]
\label{def:gm_class}
Let $m \geq 0$. The class $\mathcal{C}(m)$ holds the laws of the random
variables
\[
  S = \delta + \sum_{i=1}^{n} c_i\, Y_i, \qquad
  Y_i \sim \mathrm{Gamma}(a_i, 1) \text{ independent}, \quad
  c_i \in (0, 1], \quad a_i > 0, \quad \delta \geq 0,
\]
with $n \geq 0$ and mean $\delta + \sum_i c_i a_i = m$. The term $c\,Y$ with
$Y \sim \mathrm{Gamma}(m'/c, 1)$ has mean $m'$; as $c \to 0$ it tends to the
constant $m'$.
\end{definition}

\begin{lemma}[True states belong to the class]
\label{lem:gm_class}
Let $D_k$ be the damage of a Gamma cell on day $k$ of a schedule, under either
law. Then the law of $\beta^* D_k$ belongs to $\mathcal{C}(\beta^*\,\mathbb{E}D_k)$.
\end{lemma}

\begin{proof}
Since the last replacement (or day 0), the damage is
$D_k = \delta + \sum_t w_t X_t$: the $X_t$ are the independent Gamma terms added
since then, $\delta \geq 0$ is the weighted sum of the known constants, and each
weight $w_t \in (0, 1]$ is a product of factors $1-\rho$. A mission adds a term
with weight 1. $\ARDinf$ multiplies every weight, constants included, by
$1-\rho$. $\ARDone$ multiplies by $1-\rho$ the weights of the terms and constants
added since the last intervention, and keeps the others. A replacement starts a
new sum with the replacement draw. When the starting anchor $g^{\mu}_0$ is below
$\mu_0$, the initial damage is read as two independent terms of rate $\beta_0$,
with means $g^{\mu}_0$ and $\mu_0 - g^{\mu}_0$, whose sum has the law
$\mathrm{Gamma}(\beta_0\mu_0, \beta_0)$; a known initial damage is read as the
two constants $g^{\mu}_0$ and $\mu_0 - g^{\mu}_0$. Write $X_t = Y_t/\beta_t$ with
$Y_t \sim \mathrm{Gamma}(a_t, 1)$ and $\beta_t \geq \beta^*$ by
Eq.~\eqref{eq:gm_betastar}. Then
\[
  \beta^* D_k = \beta^*\delta + \sum_t \frac{w_t\,\beta^*}{\beta_t}\, Y_t,
  \qquad \frac{w_t\,\beta^*}{\beta_t} \in (0, 1].
\]
The $Y_t$ are independent by Assumption G1. Therefore the law of $\beta^* D_k$
is in $\mathcal{C}(\beta^*\,\mathbb{E}D_k)$.
\end{proof}

\subsubsection{The largest tail at a given mean}

For $0 < m < t$, let $W(m, t)$ be the supremum of $\Pr(S > t)$ over the random
variables $S$ whose law is in $\mathcal{C}(m)$.
$Q(a, x) = \Pr(\mathrm{Gamma}(a, 1) > x)$ is the regularised upper incomplete
gamma function (\texttt{scipy.special.gammaincc(a, x)}). The one-term member
$c\,\mathrm{Gamma}(m/c, 1)$ has the tail $Q(m/c,\, t/c)$ at $t$.

A member of $\mathcal{C}(m)$ is described by its atoms: the pairs
$(m_i, c_i)$ of a term mean $m_i = c_i a_i$ and a scale $c_i$, and the constant
$\delta$, read as an atom at scale $c = 0$. Terms with equal scales merge into
one term, because a sum of independent Gamma variables of one rate is a Gamma
variable.

\begin{lemma}[First-order rates]
\label{lem:gm_rates}
Let $S$ be a member with at least one term, with density $f_S$. Since
$\delta \leq m < t$, $f_S$ is continuous at $t$. For $c \in (0, 1]$ let
\[
  \varphi_S(c) = \frac{1}{c}\int_0^\infty \Pr(t - y < S \leq t)\,
    \frac{e^{-y/c}}{y}\, dy,
  \qquad \varphi_S(0) = f_S(t).
\]
(a) Adding mean $\eta$ at scale $c$ (a term $c\,\mathrm{Gamma}(\eta/c, 1)$,
or the constant $\eta$ when $c = 0$) raises $\Pr(S > t)$ at the rate
$\varphi_S(c)$ as $\eta \to 0$; removing mean from an atom at scale $c$ lowers
it at the same rate. (b) Moving the scale of an atom $(m_i, c_i)$,
$c_i \in (0,1)$, at fixed mean $m_i$ changes $\Pr(S > t)$ at the rate
$m_i\, \partial\varphi_S/\partial c\,(c_i)$.
\end{lemma}

\begin{proof}
$S$ is infinitely divisible, with drift $\delta$ and Lévy density
$\ell(y) = \sum_i (m_i/c_i)\, e^{-y/c_i}/y$ on $y > 0$. A change of the
parameters at the rates $(\dot\delta, \dot\ell)$ changes the Laplace transform
$\mathbb{E}e^{-\sigma S} = \exp\bigl(-\sigma\delta - \int_0^\infty (1 - e^{-\sigma y})\,
\ell(y)\, dy\bigr)$ at the rate
$\int_0^\infty \mathbb{E}\bigl[e^{-\sigma(S+y)} - e^{-\sigma S}\bigr]\,\dot\ell(y)\, dy
- \sigma\dot\delta\,\mathbb{E}e^{-\sigma S}$. This is the transform of the rate
$\int_0^\infty \Pr(t - y < S \leq t)\,\dot\ell(y)\, dy + \dot\delta\, f_S(t)$ of
$\Pr(S > t)$. Case (a) is $\dot\ell(y) = c^{-1}e^{-y/c}/y$, or $\dot\delta = 1$
when $c = 0$. Case (b) is
$\dot\ell(y) = m_i\,\partial_c\bigl(c^{-1}e^{-y/c}/y\bigr)$ at $c = c_i$.
\end{proof}

\begin{theorem}[One term is the worst case]
\label{thm:gm_oneterm}
Let $0 < m < t$. Then
\begin{equation}
  W(m, t) = \max_{c \in (0, 1]} Q\!\left(\frac{m}{c},\, \frac{t}{c}\right).
  \label{eq:gm_oneterm}
\end{equation}
\end{theorem}

\begin{proof}
Let $\mathcal{C}_K(m)$ be the members with at most $K$ atoms, the constant
included, and let $K \geq 3$.
\begin{enumerate}[leftmargin=1.5em]
  \item \emph{A maximiser exists.} $\mathcal{C}_K(m)$ is the image of the
        compact set of atoms $(m_i, c_i)$ with $m_i \geq 0$, $\sum_i m_i = m$,
        $c_i \in [0, 1]$. The law of $S$ depends continuously on the atoms, and
        only the constant member $S = m$ has an atom, at $m \neq t$. So
        $\Pr(S > t)$ is continuous and attains its maximum at some $S^*$. Since
        $\Pr(m > t) = 0 < Q(m, t)$, $S^*$ has at least one term. Every member
        has finitely many atoms, so
        $W(m, t) = \sup_K \max_{\mathcal{C}_K(m)} \Pr(S > t)$.
  \item \emph{First-order conditions.} Write $\varphi^* = \varphi_{S^*}$. By
        Lemma~\ref{lem:gm_rates}: (i) moving mean between two atoms of positive
        mean must not raise the tail, so $\varphi^*$ takes one value $\kappa^*$
        at every atom of positive mean; (ii) at an atom with
        $c \in (0, 1)$, $\partial\varphi^*/\partial c = 0$; (iii) if $S^*$ has
        fewer than $K$ atoms of positive mean, a new atom can be opened at any
        $c \in [0, 1]$, so $\varphi^*(c) \leq \kappa^*$ on $[0, 1]$.
  \item \emph{A single-peak function.} Put $\lambda = 1/c$ and
        $h(u) = \Pr(t - u < S^* \leq t)/u$. Then
        $\varphi^*(\lambda) = \lambda\int_0^\infty h(u)\, e^{-\lambda u}\, du$,
        and $\varphi^*(\infty) = h(0^+) = f_{S^*}(t)$ is the value at the
        constant atom. $S^*$ is a constant plus a sum of independent Gamma
        variables, so its density is unimodal
        \cite[Theorem~4, p.~11]{roosta2015}, proved in
        \cite[proof of Theorem~4, pp.~194--195]{szekely2003}; the constant
        shifts the mode only. Then $g(v) = f_{S^*}(t - v)$ for $v < t$ and
        $g(v) = 0$ for $v \geq t$ is unimodal on $[0, \infty)$, and so is its
        running average $h(u) = u^{-1}\int_0^u g$: before the mode of $g$,
        $g \geq h$; after it, $u\,(g(u) - h(u)) = u\,g(u) - \int_0^u g$ is
        nonincreasing. Since $h' = (g - h)/u$, $h$ rises and then falls.
  \item \emph{Zero count.} For every constant $\kappa'$, $h - \kappa'$ changes
        sign at most twice. The Laplace transform of a function with at most
        $n$ sign changes has at most $n$ zeros in $\lambda > 0$, counted with
        multiplicity: at a sign change $u_1$, multiply the transform by
        $e^{\lambda u_1}$; its derivative in $\lambda$ is the transform of a
        function with one sign change fewer; Rolle's theorem and induction
        conclude. So $\varphi^* - \kappa' = \lambda\int_0^\infty (h - \kappa')\,
        e^{-\lambda u}\, du$ has at most two zeros in $\lambda > 0$. Also
        $\varphi^*(\lambda) \to 0$ as $\lambda \to 0^+$. Hence $\varphi^*$
        cannot rise above a level $\kappa' > 0$, fall below it and rise above it
        again: with $\varphi^*(0^+) = 0$ that would give three zeros. So
        $\varphi^*$ is quasi-concave. It is analytic in $\lambda > 0$ and not
        constant, so its maximum over $\lambda \in [1, \infty]$ is attained at
        one point only: two maximisers would make $\varphi^*$ constant between
        them.
  \item \emph{One term.} By (i) and (ii), an atom with $c \in (0, 1)$ is a zero
        of $\varphi^* - \kappa^*$ of multiplicity at least 2, and an atom at
        $c = 1$ a zero of multiplicity at least 1; the constant atom is at
        $\lambda = \infty$. By step~4, $S^*$ has at most two atoms of positive
        mean. As $K \geq 3$, (iii) applies, so $\kappa^*$ is the maximum of
        $\varphi^*$ over $[1, \infty]$, attained at every atom of positive mean.
        By step~4 there is only one such atom, and by step~1 it is not the
        constant. So $S^* = c^*\,\mathrm{Gamma}(m/c^*, 1)$ with
        $c^* \in (0, 1]$.
\end{enumerate}
Therefore $\max_{\mathcal{C}_K(m)} \Pr(S > t) = \max_{c \in (0,1]} Q(m/c, t/c)$
for every $K \geq 3$, and Eq.~\eqref{eq:gm_oneterm} follows from step~1.
\end{proof}

\begin{lemma}[Monotone in the mean]
\label{lem:gm_monotone}
$W(m, t)$ is nondecreasing in $m \in (0, t)$.
\end{lemma}

\begin{proof}
Let $m < m' < t$ and let the law of $S$ be in $\mathcal{C}(m)$. Then the law of
$S + (m' - m)$ is in $\mathcal{C}(m')$, and
$\Pr(S + m' - m > t) \geq \Pr(S > t)$. Therefore $W(m, t) \leq W(m', t)$.
\end{proof}

\subsubsection{Reliability row}

\textbf{The cap.} For $A \in [0, t)$ let
$w(A) = \max_{c \in (0, 1]} Q(A/c,\, t/c)$; it is continuous and nondecreasing
in $A$, because $Q(a, x)$ increases in $a$. The \emph{sharp cap}
$A_1(t, \varepsilon)$ is the largest $A \in [0, t)$ with
$w(A) \leq \varepsilon$, computed as described below. The row, on every day
$k = 1, \ldots, T$, is
\begin{equation}
  \mu_k \;\leq\; \frac{A_1(\beta^*\tau,\, \varepsilon)}{\beta^*}.
  \label{eq:gm_row}
\end{equation}
The row is linear: it is built in its exact form, and a value of
\texttt{reliability\_impl} at the cell is ignored (Section~\ref{sec:linear_forms}).

\begin{corollary}[The row is safe]
\label{cor:gm_safe}
Let the tracked mean satisfy Eq.~\eqref{eq:gm_row} on day $k$. Then
$\Pr(D_k > \tau) \leq \varepsilon$.
\end{corollary}

\begin{proof}
Let $m = \beta^*\,\mathbb{E}D_k$ and $A = A_1(\beta^*\tau, \varepsilon)$. The
tracked mean bounds the true mean from above (Section~\ref{sec:encoding}), so
$m \leq \beta^*\mu_k \leq A < t$. If $m = 0$, then $D_k = 0$. Otherwise
\[
  \Pr(D_k > \tau) = \Pr(\beta^* D_k > t)
  \leq W(m, t) \leq W(A, t) = w(A) \leq \varepsilon,
\]
by Lemma~\ref{lem:gm_class}, Lemma~\ref{lem:gm_monotone},
Theorem~\ref{thm:gm_oneterm} and the choice of $A$.
\end{proof}

\textbf{Computation of $A_1$} (offline, once per cell, before the model is built):
\begin{enumerate}[leftmargin=1.5em, nosep]
  \item $t = \beta^*\tau$. $\hat A$ is the root of $Q(\hat A, t) = \varepsilon$,
        found by bisection on $\hat A$ ($Q$ increases in its first argument;
        \texttt{gammaincinv} solves for the second argument and does not apply).
  \item $w(A)$: the function $c \mapsto Q(A/c, t/c)$ tends to 0 as $c \to 0$
        and has a single peak on $(0, 1]$ (checked numerically for $t$ from
        $10^{-3}$ to $100$ and $A/t$ from $0.05$ to $0.999$). Its maximum is found on a
        logarithmic grid of at least 400 values of $c$ in
        $[\min(10^{-4}, t/100),\, 1]$, then refined by a golden-section search
        between the neighbours of the best grid point. If the best grid point is
        the smallest value of $c$, the grid is extended downwards.
  \item If $\hat A < t$ and $w(\hat A) \leq \varepsilon$, then
        $A_1 = \hat A$. Otherwise $A_1$ is found by bisection on $A$ in
        $[\varepsilon t,\, \min(\hat A, t))$, to a relative tolerance of
        $10^{-9}$. The lower end satisfies $w(\varepsilon t) \leq \varepsilon$,
        by Markov's inequality $Q(a, x) \leq a/x$. The bisection returns the
        lower end of its last interval, so $w(A_1) \leq \varepsilon$ holds at
        the returned value.
\end{enumerate}

\emph{Remark (the sharp cap and the exact one-rate cap).} $\hat A$ is the exact
cap of a single Gamma law of rate $\beta^*$, and $A_1 \leq \hat A$, because
$c = 1$ is in the maximum. Let $t_1(\varepsilon)$ be the value of $t$ at which
$\partial_c Q(\hat A/c, t/c)$ at $c = 1$ changes sign: $t_1 = 0.436$, $0.448$
and $0.614$ for $\varepsilon = 0.001$, $0.01$ and $0.1$. For $t \geq t_1$,
$A_1 = \hat A$: on a grid of $t$ in $[t_1, 1000]$, no $c < 1$ gives a larger tail
than $c = 1$. Below $t_1$, $A_1 < \hat A$. For $\varepsilon = 0.01$:
\begin{center}
\begin{tabular}{@{}lccccc@{}}
\toprule
$t = \beta^*\tau$ & 0.01 & 0.1 & 0.3 & 1 & 10 \\
\midrule
$A_1/\hat A$ & 0.142 & 0.642 & 0.961 & 1 & 1 \\
\bottomrule
\end{tabular}
\end{center}
When every term has the same rate and no repair occurs, the true state is a single
Gamma law and $\hat A$ is exact for it. With several rates, or with repairs, the true state is a
weighted sum, and below $t_1$ the exact one-rate cap is not safe. Example:
$\varepsilon = 0.01$ and $\beta^*\tau = 0.1$, so $\hat A = 0.0055012$. A new
part flies one mission of mean $\hat A/\beta^*$ and rate $4.56\,\beta^*$. Its
mean meets the exact one-rate cap, but its tail is $1.565\,\varepsilon$. The
sharp cap is $A_1 = 0.642\,\hat A$ there.

\emph{Remark (below the threshold).} With $r = A/t$ and $s = t/c$,
$Q(A/c, t/c) = Q(r s, s)$, and $c \in (0, 1]$ means $s \geq t$. For
$t \leq t_1(\varepsilon)$ the peak of $s \mapsto Q(rs, s)$ at the cap lies at
$s \geq t$, so the maximum does not depend on $t$, and
$A_1(t, \varepsilon) = r^*(\varepsilon)\, t$. The mean cap is then
$A_1/\beta^* = r^*(\varepsilon)\,\tau$, whatever the rates:
$r^* = 0.00355$, $0.0353$ and $0.331$ for $\varepsilon = 0.001$, $0.01$ and
$0.1$.

\emph{Remark (the choice of $\beta^*$).} A larger $t$ shrinks the range
$s \geq t$, so $A_1(t, \varepsilon)/t$ is nondecreasing in $t$, and the mean cap
$A_1(\beta^*\tau, \varepsilon)/\beta^* = \tau\, A_1(t, \varepsilon)/t$ is
nondecreasing in $\beta^*$. Lemma~\ref{lem:gm_class} holds for every
$\beta^*$ at most the smallest rate. So the smallest rate gives the largest safe
cap of this form.

\emph{Remark (conservatism).} The class $\mathcal{C}(m)$ also contains laws that
no schedule produces. Below $t_1$ the row can therefore reject some safe
schedules. The row keeps each mission's own mean; $\beta^*$ only bounds how
dispersed each term may be.

\textbf{Facts used by the encoding and loop arguments.}
(M1) The day map~\eqref{eq:gm_map} is nondecreasing in $(\mu, g^{\mu})$.
(M2) The true mean follows the map~\eqref{eq:gm_map} exactly.
(M3) If the tracked mean satisfies Eq.~\eqref{eq:gm_row} and the true mean is
not larger, then $\Pr(D_k > \tau) \leq \varepsilon$
(Corollary~\ref{cor:gm_safe}). The loop row of Section~\ref{sec:loop} covers
$\mu$ and, under $\ARDone$, $g^{\mu}$.

% ----------------------------------------------------------------
\subsection{Rainflow Model}
\label{sec:dyn_rainflow}
% ----------------------------------------------------------------

This section applies to the cells $(i,\ell)$ with \texttt{model} $=$
\texttt{rainflow}. Inside it, the subscript $i\ell$ is dropped from the per-cell
parameters and from the states; $k$ is the day.

The rainflow model is \emph{distribution-free}. The law of a mission increment
is measured from recorded load histories by rainflow counting (Palmgren--Miner
damage), and it enters the model only through summary statistics: its mean, and,
as the bound requires, its variance, the width of its range, or its cumulant
generating function (CGF) at a fixed point $s$. No parametric family is assumed.
The reliability requirement $\Pr(D_k > \tau) \leq \varepsilon$ is certified by
one of five tail bounds, chosen per cell by \texttt{bound\_method}:
\texttt{markov}, \texttt{cantelli}, \texttt{hoeffding}, \texttt{bernstein},
\texttt{chernoff}.

\textbf{Modelling assumptions (normative).}
\begin{itemize}[leftmargin=1.5em]
  \item \textbf{A0 (non-negative damage).} Every increment, the initial damage
        and every replacement draw are non-negative.
  \item \textbf{A1 (independence).} The mission increments are mutually
        independent, and independent of the initial damage and of every
        replacement draw. Each replacement installs a new draw, independent of
        everything else. Every bound except Markov needs A1.
  \item \textbf{A2 (bounded ranges; Hoeffding and Bernstein).} Each mission
        increment lies in an interval of width $b^{\inc}_{jk}$ (input
        \texttt{support}). The initial damage is a known part plus independent
        terms, each in an interval; $R_0$ is the sum of their squared widths
        (input \texttt{R\_0}) and $b^0$ the largest width (input
        \texttt{support\_0}). The replacement draw is described in the same way by
        $R^{\new}$ and $b^{\new}$ (inputs \texttt{R\_new}, \texttt{support\_new}).
        A known damage has $R = b = 0$, the defaults.
  \item \textbf{A3 (repetition).} The row of the cell's bound holds on every day
        $k = 1, \ldots, T$, and after day $T$ the operating phase repeats with
        period $H_2$, schedule and increment profile included. This is the
        setting of the loop argument of Section~\ref{sec:loop}.
\end{itemize}

\textbf{Width of a range.} The input \texttt{support} of an increment $X$ is the
width $b$ of an interval $[l, l + b]$ that contains $X$ almost surely. Hoeffding's
inequality uses the widths. Bernstein's inequality needs an upper bound of the
centred term, and the width is one: $X - \mathbb{E}X \leq l + b - l = b$, since
$\mathbb{E}X \geq l$. A width $b$ is consistent with a variance $v$ only if
$v \leq b^2/4$ (Popoviciu's inequality:
$\mathrm{Var}\,X \leq \mathbb{E}(X - l - b/2)^2 \leq (b/2)^2$); the check is in
Section~\ref{sec:consistency_checks}.

\textbf{Tracked states per bound.} Every cell tracks the mean $\mu$. Each bound
adds at most one state, and under $\ARDone$ one anchor per tracked state
($g^{\mu}$ and $g^{v}$, $g^{R}$ or $g^{K}$). Both laws are accepted for every
bound.

\begingroup\small
\begin{longtable}{@{}lllll@{}}
\toprule
\texttt{bound\_method} & States & Increment inputs & Initial; replacement inputs & Forms \\
\midrule
\texttt{markov} & $\mu$ & \texttt{mu} & -- & \texttt{exact} \\
\texttt{cantelli} & $\mu$, $v$ & \texttt{mu}, \texttt{v}
  & \texttt{v\_0}; \texttt{replacement\_v} & all three \\
\texttt{hoeffding} & $\mu$, $R$ & \texttt{mu}, \texttt{support}
  & \texttt{R\_0}; \texttt{R\_new} & all three \\
\texttt{bernstein} & $\mu$, $v$ & \texttt{mu}, \texttt{v}, \texttt{support}
  & \texttt{v\_0}, \texttt{support\_0}; & all three \\
 & & & \texttt{replacement\_v}, \texttt{support\_new} & \\
\texttt{chernoff} & $\mu$, $K$ & \texttt{mu}, \texttt{cgf}, \texttt{s\_chernoff}
  & \texttt{cgf\_0}; \texttt{cgf\_new} & \texttt{exact} \\
\bottomrule
\end{longtable}
\endgroup

The column ``Forms'' lists the forms in which the row can be built; ``all
three'' means \texttt{exact}, \texttt{tangent} and \texttt{pwl}. A linear row is
built in its exact form, and a value of \texttt{reliability\_impl} at its cell is
ignored.
$R$ is the \emph{Hoeffding budget}, the sum of the squared widths of the terms of
the damage. $K$ is the CGF of the damage at the cell's parameter
$s = \texttt{s\_chernoff} > 0$. The day increments are
$\delta\mu_k = \sum_{j=1}^M x_{ijk}\,\mu^{\inc}_{jk}$,
$\delta v_k = \sum_{j} x_{ijk}\,v^{\inc}_{jk}$,
$\delta R_k = \sum_{j} x_{ijk}\,(b^{\inc}_{jk})^2$ and
$\delta K_k = \sum_{j} x_{ijk}\,K^{\inc}_{jk}$, with
$K^{\inc}_{jk} = \ln\mathbb{E}\,e^{s X_{jk}}$ the CGF of the increment at $s$
(input \texttt{cgf}).

\subsubsection{Day maps}

Between replacements the damage is $D_k = \sum_t w_t X_t$: the $X_t$ are the
known part, the independent terms of the initial damage or of the replacement
draw, and the mission increments; each weight $w_t \in (0, 1]$ is a product of
factors $1 - \rho$. Under $\ARDone$, write $D^- = G + \Delta$, with $G$ the damage
just after the last intervention and $\Delta$ the weighted increments added
since. $G$ and $\Delta$ are independent by A1, and a repair gives
$D^+ = G + (1-\rho)\,\Delta$. Before the first intervention, $G$ is the part of
the initial damage that the starting anchors describe
(Section~\ref{sec:repair_laws}).

For a damage $D = \sum_t w_t X_t$, write $m(D) = \mathbb{E}D$,
$V(D) = \mathrm{Var}\,D$, $R(D) = \sum_t w_t^2 r_t$ with $r_t$ the squared width
of $X_t$, and $K(D) = \ln\mathbb{E}\,e^{sD}$.

\begin{proposition}[Day maps of the statistics]
\label{prop:rf_maps}
Let $\theta \in \{m, V, R\}$ with factor $a_\theta = 1-\rho$ for $m$ and
$a_\theta = (1-\rho)^2$ for $V$ and $R$. Under A1:
\begin{enumerate}[leftmargin=1.5em, nosep]
  \item a mission day adds the increment's statistic:
        $\theta(D + X) = \theta(D) + \theta(X)$ and $K(D + X) = K(D) + K(X)$;
  \item $\ARDinf$: $\theta\bigl((1-\rho)D\bigr) = a_\theta\,\theta(D)$ and
        $K\bigl((1-\rho)D\bigr) \leq (1-\rho)\,K(D)$;
  \item $\ARDone$: $\theta(D^+) = a_\theta\,\theta(D^-) + (1 - a_\theta)\,\theta(G)$
        and $K(D^+) \leq (1-\rho)\,K(D^-) + \rho\,K(G)$.
\end{enumerate}
\end{proposition}

\begin{proof}
Item 1: the mean and $R$ are additive by definition, the variance and the CGF by
independence. Item 2: $m$ and $V$ scale by $1-\rho$ and $(1-\rho)^2$; every
weight is multiplied by $1-\rho$, so $R$ scales by $(1-\rho)^2$. A CGF is convex
with value 0 at 0, so
$K\bigl((1-\rho)D\bigr) = K_D\bigl((1-\rho)s\bigr) \leq (1-\rho)\,K_D(s) + \rho\,K_D(0)
= (1-\rho)\,K(D)$, with $K_D$ the CGF of $D$ as a function of its argument.
Item 3: by independence, $\theta(D^-) = \theta(G) + \theta(\Delta)$; for $R$ the
weights of the terms of $G$ are unchanged and those of $\Delta$ are multiplied by
$1-\rho$. Hence
\[
  \theta(D^+) = \theta(G) + a_\theta\,\theta(\Delta)
  = \theta(G) + a_\theta\bigl(\theta(D^-) - \theta(G)\bigr)
  = a_\theta\,\theta(D^-) + (1 - a_\theta)\,\theta(G),
\]
and, with item 2 applied to $\Delta$,
\[
  K(D^+) = K(G) + K_\Delta\bigl((1-\rho)s\bigr)
  \leq K(G) + (1-\rho)\bigl(K(D^-) - K(G)\bigr)
  = (1-\rho)\,K(D^-) + \rho\,K(G).
\]
\end{proof}

So each tracked state follows the generic map of Section~\ref{sec:encoding} with
the factor $a = 1-\rho$ for $\mu$ and $K$, and $a = (1-\rho)^2$ for $v$ and $R$.
Its anchor is its value just after the last intervention. The maps of $\mu$,
$v$ and $R$ are exact; the map of $K$ is an upper bound, which is what the
domination argument of Section~\ref{sec:encoding} needs.

\subsubsection{Starting and replacement values}

The initial and replacement values of $\mu$ and $v$ come from the inputs
(Section~\ref{sec:replacement}). Those of $R$ and $K$ are:
\begin{itemize}[leftmargin=1.5em, nosep]
  \item $R_0$ from \texttt{R\_0} and $R^{\new}$ from \texttt{R\_new}, default 0;
  \item $K_0$ from \texttt{cgf\_0}, default $s\,\mu_0$, and $K^{\new}$ from
        \texttt{cgf\_new}, default $s\,\mu^{\new}$.
\end{itemize}
With \texttt{replacement\_as\_new} all replacement values are 0
(Section~\ref{sec:replacement}).

\emph{Why the defaults.} By Jensen's inequality,
$K(D) = \ln\mathbb{E}\,e^{sD} \geq s\,\mathbb{E}D$, with equality exactly when $D$
is almost surely constant. So $s\mu_0$ is the exact CGF of a known initial
damage $\mu_0$, and a lower bound for a random one; a random initial damage needs
its own CGF in \texttt{cgf\_0}. The input checks of
Section~\ref{sec:consistency_checks} require $\texttt{cgf} \geq s\,\mu^{\inc}$ for
every increment, $K_0 \geq s\mu_0$ and $K^{\new} \geq s\mu^{\new}$. A known damage
has width 0, so $R_0 = 0$ is exact for it, and a random one needs \texttt{R\_0}.
The same holds for the replacement draw. A starting value of $K$ below $s\mu_0$ makes the row optimistic. Example:
$\tau = 1$, $\varepsilon = 0.1$, $s = 5$, a known initial damage $\mu_0 = 0.6$,
and one mission that adds 0 or 0.5 with probability $1/2$ each, so
$K^{\inc} = \ln(0.5 + 0.5\,e^{2.5}) = 1.886$. Then
$\Pr(D > \tau) = 0.5$. The row with $K_0 = s\mu_0 = 3$ gives
$4.886 - 5 > \ln 0.1$ and rejects the plan; a start at $K_0 = 0$ would give
$1.886 - 5 = -3.11 \leq \ln 0.1 = -2.30$ and accept it.

\subsubsection{Rows of the dynamics}

The rows are the generic rows of Section~\ref{sec:encoding}, built for $\mu$
and for the one extra state of the bound, with the data below.
$V^{\mathrm{C}} = \varepsilon\tau^2/(1-\varepsilon)$ and
$V^{\mathrm{B}} = \tau(\tau - 2c)/(2\kappa)$ are the variance ranges of the
Cantelli and Bernstein rows, with $\kappa = \ln(1/\varepsilon)$ and $c$ from
Eq.~\eqref{eq:rf_bsupport}. Each natural range is the largest value that the
bound's exact row admits at $\mu \geq 0$. The big-M rows need every initial
value, starting anchor and replacement value inside the natural range of its
state; the input checks of Section~\ref{sec:consistency_checks} require it.

\begingroup\small
\begin{longtable}{@{}llllll@{}}
\toprule
State & Day increment & Factor $a$ & Natural range & Initial value & Replacement value \\
\midrule
$\mu$; anchor $g^{\mu}$ & $\delta\mu_k$ & $1-\rho$ & $\tau$ & $\mu_0$ & $\mu^{\new}$ \\
$v$; anchor $g^{v}$ & $\delta v_k$ & $(1-\rho)^2$ & $V^{\mathrm{C}}$ or $V^{\mathrm{B}}$
  & $v_0$ & $v^{\new}$ \\
$R$; anchor $g^{R}$ & $\delta R_k$ & $(1-\rho)^2$ & $2\tau^2/\kappa$ & $R_0$ & $R^{\new}$ \\
$K$; anchor $g^{K}$ & $\delta K_k$ & $1-\rho$ & $s\tau - \kappa$ & $K_0$ & $K^{\new}$ \\
\bottomrule
\end{longtable}
\endgroup

The anchors exist only under $\ARDone$, and the extra state only for its bound:
$v$ for Cantelli and Bernstein, $R$ for Hoeffding, $K$ for Chernoff.

\subsubsection{Reliability rows}

Write $d_k = \tau - \mu_k$. Let $(m, V, R^{\mathrm{t}}, K^{\mathrm{t}})$ be the
true statistics of $D_k$. By Proposition~\ref{prop:rf_maps} and
Section~\ref{sec:encoding}, each is at most the tracked state of the same name.
Each row below is imposed on every day $k = 1, \ldots, T$.

\textbf{Markov} (\texttt{markov}). By A0 and Markov's inequality,
$\Pr(D_k > \tau) \leq m/\tau \leq \mu_k/\tau$. The row is
\begin{equation}
  \mu_k \leq \varepsilon\,\tau.
  \label{eq:rf_markov}
\end{equation}
It uses the mean only and does not need A1 or A2.

\textbf{Cantelli} (\texttt{cantelli}). Cantelli's inequality gives
$\Pr(D - m \geq a) \leq V/(V + a^2)$ for $a > 0$. The row is
\begin{equation}
  (1 - \varepsilon)\, v_k \leq \varepsilon\,(\tau - \mu_k)^2, \qquad \mu_k \leq \tau,
  \label{eq:rf_cantelli}
\end{equation}
equivalently $\mu_k + \kappa_C\sqrt{v_k} \leq \tau$ with
$\kappa_C = \sqrt{(1-\varepsilon)/\varepsilon}$. \emph{Validity.} If $d_k > 0$,
\[
  \Pr(D_k > \tau) \leq \Pr(D_k - m \geq \tau - m) \leq \frac{V}{V + (\tau - m)^2}
  \leq \frac{v_k}{v_k + d_k^2} \leq \varepsilon,
\]
because $v/(v + d^2)$ increases in $v$ and decreases in $d$, and $\tau - m \geq d_k$.
If $d_k = 0$, the row gives $v_k = 0$, so $V = 0$ and $D_k = m \leq \tau$.

\textbf{Hoeffding} (\texttt{hoeffding}). $D_k - m$ is a sum of independent
centred terms $w_t(X_t - \mathbb{E}X_t)$ with ranges of widths $w_t\sqrt{r_t}$,
and the sum of their squares is $R^{\mathrm{t}}$. Hoeffding's inequality gives
$\Pr(D_k - m \geq a) \leq \exp(-2a^2/R^{\mathrm{t}})$, and
$\exp(-2d^2/R) \leq \varepsilon$ exactly when $d^2 \geq (\kappa/2)\,R$. The row
is
\begin{equation}
  (\tau - \mu_k)^2 \geq \frac{\kappa}{2}\, R_k, \qquad \mu_k \leq \tau.
  \label{eq:rf_hoeffding}
\end{equation}
\emph{Validity.} If $R^{\mathrm{t}} = 0$, every term is constant and
$D_k = m \leq \tau$; this covers $d_k = 0$, where the row gives $R_k = 0$.
Otherwise $d_k > 0$, and
$\Pr(D_k > \tau) \leq \exp\bigl(-2(\tau - m)^2/R^{\mathrm{t}}\bigr)
\leq \exp(-2d_k^2/R_k) \leq \varepsilon$, because $\exp(-2d^2/R)$ increases in
$R$ and decreases in $d$, and $\tau - m \geq d_k$.

\textbf{Bernstein} (\texttt{bernstein}). The \emph{support constant} of the cell
is
\begin{equation}
  b = \max\Bigl(b^0,\; b^{\new},\; \max_{j,k} b^{\inc}_{jk}\Bigr), \qquad
  c = \frac{\kappa\, b}{3},
  \label{eq:rf_bsupport}
\end{equation}
over the missions $j$ and the days $k$ of both phases. It is a constant, not a
decision variable. Every weighted term satisfies
$w_t(X_t - \mathbb{E}X_t) \leq w_t\, b_t \leq b$, whatever the schedule, because
$w_t \leq 1$. Bernstein's inequality for independent terms bounded above by $b$
gives $\Pr(D - m \geq a) \leq \exp\bigl(-a^2/(2(V + b\,a/3))\bigr)$ for $a > 0$;
the right side is at most $\varepsilon$ exactly when
$a^2 - 2c\,a - 2\kappa V \geq 0$. The row is
\begin{equation}
  (\tau - c - \mu_k)^2 \geq c^2 + 2\kappa\, v_k, \qquad \mu_k \leq \tau - c,
  \label{eq:rf_bernstein}
\end{equation}
equivalently $\mu_k + c + \sqrt{c^2 + 2\kappa v_k} \leq \tau$. With $v_k \geq 0$,
the two rows give $d_k \geq 2c$.
\emph{Validity.} $a^2/(2(V + b\,a/3))$ increases in $a > 0$ and decreases in $V$,
and $\tau - m \geq d_k \geq 2c$; so $\Pr(D_k > \tau) \leq \varepsilon$ when
$b > 0$. When $b = 0$, every term is constant and $D_k = m \leq \tau$.

\emph{Remark (the side row).} The quadratic row alone also admits the point
$d_k = 0$, $v_k = 0$: a known damage equal to $\tau$. That point is safe, but it is
cut off from the rest of the feasible set. The side row
$\mu_k \leq \tau - c$ removes it and keeps the branch
$d_k \geq c + \sqrt{c^2 + 2\kappa v_k}$. At $\mu = v = 0$ the row needs
$\tau \geq 2c$; a positive variance needs $\tau > 2\kappa b/3$, the input check
of Section~\ref{sec:consistency_checks}.

\textbf{Chernoff} (\texttt{chernoff}). For $s > 0$, Markov's inequality applied
to $e^{sD}$ gives $\Pr(D_k \geq \tau) \leq \exp(K^{\mathrm{t}} - s\tau)$. The row
is
\begin{equation}
  K_k - s\,\tau \leq \ln\varepsilon.
  \label{eq:rf_chernoff}
\end{equation}
\emph{Validity.} $K^{\mathrm{t}} \leq K_k$, so
$\Pr(D_k > \tau) \leq \exp(K_k - s\tau) \leq \varepsilon$. The row bounds $K$ by
$s\tau - \kappa$, its natural range; $K \geq 0$, because
$K(X) \geq s\,\mathbb{E}X \geq 0$ by Jensen's inequality and A0. A known zero
damage satisfies the row only if $s\tau \geq \kappa$ (input check
$s > \ln(1/\varepsilon)/\tau$, Section~\ref{sec:consistency_checks}). A
replacement day satisfies it only if $K^{\new} \leq s\tau - \kappa$; the input
checks require every replacement state to satisfy the cell's row.

\textbf{Linear forms.} Markov and Chernoff rows are linear: they are built in
their exact form, and a value of \texttt{reliability\_impl} at the cell is
ignored. The Cantelli, Hoeffding and Bernstein rows have the form
$\chi_k \leq f(\mu_k)$ with $f$ convex in $\mu$, and accept the forms
\texttt{exact}, \texttt{tangent} and \texttt{pwl} of
Section~\ref{sec:linear_forms}. Every form keeps the side row.

\begingroup\small
\begin{longtable}{@{}lllll@{}}
\toprule
Bound & $\chi$ & $f(\mu)$ & Side row & $f(0)$ \\
\midrule
Cantelli & $v$ & $\varepsilon(\tau - \mu)^2/(1 - \varepsilon)$ & $\mu \leq \tau$
  & $\varepsilon\tau^2/(1-\varepsilon)$ \\
Hoeffding & $R$ & $2(\tau - \mu)^2/\kappa$ & $\mu \leq \tau$ & $2\tau^2/\kappa$ \\
Bernstein & $v$ & $\bigl((\tau - \mu)^2 - 2c(\tau - \mu)\bigr)/(2\kappa)$
  & $\mu \leq \tau - c$ & $\tau(\tau - 2c)/(2\kappa)$ \\
\bottomrule
\end{longtable}
\endgroup

Each $f$ is a parabola in $\mu$ with a positive leading coefficient, so it is
convex. For Bernstein with $b > 0$, the tangent point must satisfy
$\mu_r < \tau - c$, an input check of Section~\ref{sec:consistency_checks}.
At $\mu_r \geq \tau - c$ the tangent has a nonnegative
slope and lies below $f$, so at every $\mu \leq \tau - c$ it is at most
$f(\tau - c) = -c^2/(2\kappa) < 0$, and no point with $v \geq 0$ satisfies the
row.

\textbf{Optional row.} If \texttt{v\_max\_user} is given for a Cantelli or
Bernstein cell, the row $v_k \leq V^{\mathrm{user}}$ is added on every day.

\textbf{Facts used by the encoding and loop arguments.}
(M1) Every day map is nondecreasing in each state and anchor: its
coefficients $a$ and $1 - a$ lie in $[0, 1]$.
(M2) The true statistics follow the maps with equality for $\mu$, $v$ and $R$,
and with $\leq$ for $K$ (Proposition~\ref{prop:rf_maps}).
(M3) If the tracked states satisfy the row of the cell's bound and the true
statistics are not larger, then $\Pr(D_k > \tau) \leq \varepsilon$ (the validity
paragraphs above). The loop row of Section~\ref{sec:loop} covers $\mu$, the extra
state of the bound and, under $\ARDone$, their anchors.
